% ===== 原文 pp.21–23 =====
% （§4 标题、§4.1、§4.2 开头；本块 tag 4.1–4.7）

% [原文 p.21]
\section{两个超对称背景}

本节我们首先讨论超对称局域化的一般思想，遵循的是 [21] 开创、并在文集 [22] 中综述过的进路。随后我们讨论两个显式的超对称背景，一个基于 $S^1\times S^3$，另一个基于 $S^4$。关于前者与后者的更多细节，分别可以在文集 [23] 与文集 [24] 中找到。

\subsection{超对称局域化的一般心法}

% [Tutor Note: p.21, §4.1 第一段：原文把度规变分写成 $\delta\log Z=\int\delta g_{\mu\nu}T^{\mu\nu}\dd^dx$，
%  这是示意式。若取 $T^{\mu\nu}=\frac{2}{\sqrt g}\frac{\delta S}{\delta g_{\mu\nu}}$，则严格的关系是
%  $\delta\log Z=-\frac12\int\sqrt g\,T^{\mu\nu}\delta g_{\mu\nu}\dd^dx$；原式把 $\sqrt g$ 与 $\frac12$
%  吸收进积分测度，并略去了整体的符号约定。后文 (4.1)(4.4) 只用到「等于 0」，故这里的正负号不影响结论；
%  但若要真的从配分函数反解 $T^{\mu\nu}$，必须把 $\sqrt g$ 与 $\frac12$ 补回来。]

量子场论的能量–动量张量 $T_{\mu\nu}$ 描述它如何与外部度规微扰耦合。具体地说，设 $Z$ 为作为度规泛函的配分函数，则我们有 $\delta\log Z=\int\delta g_{\mu\nu}T^{\mu\nu}\dd^d x$。把这一微小变分积分起来，我们就知道量子场论在一般弯曲流形上如何行为。

在超对称理论中，能量–动量张量 $T_{\mu\nu}$ 位于一个超多重态之中，该超多重态还包含超流 $S_{\mu\alpha}$ 以及其它分量，具体情况取决于时空维数、超对称的数目以及其它更精细的性质。本综述关心的是四维 $\cN=2$ 理论。此时我们有 R 流 $J^R_\mu$ 与一个标量分量 $X$，二者都处于 $\SU(2)_R$ 的伴随表示中。

能量–动量张量 $T_{\mu\nu}$ 知道理论如何与度规 $g_{\mu\nu}$ 耦合。类似地，超流 $S_{\mu\alpha}$ 知道它如何与引力微子背景 $\psi_{\mu\alpha}$ 耦合；R 流 $J^R_\mu$ 知道它如何与 R 对称性背景 $A^R_\mu$ 耦合；标量分量 $X$ 知道它如何与标量背景 $M$ 耦合。简而言之，超多重态 $(T_{\mu\nu},S_{\mu\alpha},J^R_\mu,X,\ldots)$ 知道它如何与外部超引力背景 $(g_{\mu\nu},\psi_{\mu\alpha},A^R_\mu,M,\ldots)$ 耦合。数学文献中关于超流形有许多种定义，但从超对称场论的观点看，弯曲流形最自然的「超」版本，是一个指定了完整超引力背景的流形。

% [原文 p.22]
% （原文 p.22 的第一行是上一段末句的末尾 "specified."；中文语序无法在词间断开，故该整句留在上一页。）
这一观点最早在 [21] 中得到强调，并在文集 [22] 中有更详细的综述。

类似地，如果理论具有整体味对称群 $G$，它就有味流 $J_\mu$，后者知道理论如何与味对称性背景 $A_\mu$ 耦合。当理论是四维 $\cN=2$ 超对称理论时，$J_\mu$ 位于一个超多重态中，该超多重态包含一个标量算符 $K$，它知道如何与背景标量 $\Phi$ 耦合，其中 $K$ 与 $\Phi$ 都是 $G$ 的伴随表示。事实上，$\cN=2$ 理论的质量项不过是这一构造的特例，正如我们在 §3.1 中所回顾的。

% [Tutor Note: p.22, $\langle\delta_\xi O\rangle=0$ 的来源：设路径积分测度与作用量在等距变换 $\xi$ 下不变，
%  则 $0=\int\cD\varphi\,\delta_\xi(O\ee^{-S})=\int\cD\varphi\,(\delta_\xi O)\ee^{-S}
%  +\int\cD\varphi\,O(-\delta_\xi S)\ee^{-S}$；等距/超对称意味着 $\delta_\xi S=0$，故
%  $\langle\delta_\xi O\rangle=0$。这一步正是 (4.1) 与 (4.4) 的全部依据，原文把它当作显然跳过了。]

现在，给定一个具有等距变换 $\xi$ 的弯曲流形上的量子场论，对任意算符 $O$ 都有 $\langle\delta_\xi O\rangle=0$。类似地，给定一个超对称背景（即至少有一个超等距变换 $\epsilon$ 的超引力背景）上的超对称场论，对任意算符 $O$ 都有 $\langle\delta_\epsilon O\rangle=0$。

拉格朗日量中许多自然的玻色型算符常常可以写成 $\delta_\epsilon O$ 的形式。例如，设拉格朗日量中某个耦合 $\lambda$ 乘着一个算符 $X=\delta_\epsilon O$，那么配分函数与 $\lambda$ 无关，因为
\begin{equation}
\frac{\partial}{\partial\lambda}\log Z=\langle X\rangle=\langle\delta_\epsilon O\rangle=0.
\tag{4.1}
\end{equation}

% [Tutor Note: p.22, (4.1) 与 (4.4) 的右端在通常约定下带一个负号：若 $Z=\int\cD\varphi\,\ee^{-S}$、
%  $S\supset\lambda X$，则 $\partial_\lambda\log Z=-\langle X\rangle$；同理 $S(t)=S+t\int\dd^dx\delta_\epsilon O$
%  给出 $\partial_t\log Z(t)=-\int\dd^dx\langle\delta_\epsilon O\rangle$。原文两式都写成正号，
%  因为只用到「等于 0」，符号不影响推理；但不要把这两个正号当作可推广的恒等式抄进别处。]

还要注意，超流自身的超等距变分 $\delta_\epsilon S_{\mu\alpha}$，由能量–动量张量与超多重态中其它玻色分量线性组合给出。这意味着：度规的某种变分可以与 R 对称性背景、标量背景相应的特定变分组合起来，使得配分函数与它无关。这就使超对称背景上的配分函数对度规的具体选择不敏感。有时它只依赖于时空的拓扑或复结构。当 $\delta_\epsilon^2$ 生成背景的一个玻色型等距变换时，配分函数就只依赖于该等距变换的拓扑性质，等等。

我们还常常可以选取一个费米型算符 $O$，使得
\begin{equation}
\delta_\epsilon^2 O=0,\qquad \delta_\epsilon O\simeq\sum_\psi|\delta\psi|^2,
\tag{4.2}
\end{equation}
其中 $\psi$ 跑遍理论中所有的动力学费米子场。于是考虑形变
\begin{equation}
S\to S(t)=S+t\int\dd^d x\,\delta_\epsilon O.
\tag{4.3}
\end{equation}
由于 $\delta_\epsilon^2O=0$，形变后的拉格朗日量仍然在超等距变换 $\epsilon$ 下不变。于是对任意 $t$ 有
\begin{equation}
\frac{\partial}{\partial t}\log Z(t)=\int\dd^d x\langle\delta_\epsilon O\rangle=0.
\tag{4.4}
\end{equation}
在大 $t$ 极限下，积分局域化到满足
\begin{equation}
\delta\psi=0
\tag{4.5}
\end{equation}
的组态上，
% [原文 p.23]
以及这些组态附近的涨落。出于这个原因，我们常把 $\delta_\epsilon O$ 称作局域化项。

% [Tutor Note: p.22–23, (4.2)→(4.6) 的逻辑链：由 (4.2) 的 $\delta_\epsilon O\simeq\sum_\psi|\delta\psi|^2\ge0$，
%  形变权重 $\ee^{-t\int\delta_\epsilon O}$ 在 $t\to\infty$ 时只由 $\delta\psi=0$ 的组态（及其附近）支配；
%  又因 $Z(t)$ 与 $t$ 无关（(4.4)），可以放心取 $t\to\infty$。在局域化轨迹附近把总作用量展开，
%  涨落二次项带因子 $t$、三次以上项不带，故相互作用项相对被 $1/t$ 压低——这就是 (4.6) 里
%  只剩 $Z_{\rm classical}Z_{\rm quadr.\ fluct.}$ 的原因。注意 $(4.2)$ 的第一个条件
%  $\delta_\epsilon^2O=0$ 正是「形变后仍保持超对称」的保证：$\delta_\epsilon S(t)=t\int\delta_\epsilon^2O=0$。]

我们用空间 $\cM=\sqcup_i\cM_i$ 来参数化解 (4.5) 的空间，其中 $i$ 是区分各连通分量的某个标号。于是我们有等式
\begin{equation}
Z=\sum_i\int_{\cM_i}Z_{\text{classical}}Z_{\text{quadr.\ fluct.}},
\tag{4.6}
\end{equation}
其中 $Z_{\text{classical}}$ 是经典作用量在满足 (4.5) 的组态上取值后的指数，$Z_{\text{quadr.\ fluct.}}$ 是玻色与费米涨落的高斯积分结果。相互作用项在 $t\to\infty$ 极限下不贡献。

常常每个分量 $\cM_i$ 都是有限维的，因此配分函数 $Z$ 表现为若干显式多重积分之和。计算它仍然是一件艰巨的任务，但比原来的无穷维路径积分表达式要简单无数倍。

现在我们来讨论四维 $\cN=2$ 理论的两类特殊的超对称背景，它们的局域化已被详细做出。第一类是 $S^1\times S^3$，第二类是 $S^4$。我们不讨论下文所述事实的详细推导。幸运的是，所有细节都可以在文集 [23]（对 $S^1\times S^3$）与文集 [24]（对 $S^4$）中找到。

\subsection{$S^1\times S^3$}

首先，我们考虑 $S^1\times S^3$ 上的 $\cN=2$ 超共形理论。$\cN=2$ 超共形理论具有 $\SU(2)_R\times\U(1)_R$ 的 R 对称性。可以施行超对称局域化的一类超对称背景，由 $S^1$ 与 $S^3$ 半径之比 $\beta$、以及 $\SU(2)_R$ 与 $\U(1)_R$ 绕 $S^1$ 的两个和乐所指定。总共有三个参数，通常记作 $(p,q,t)$。注意这里的 $q$ 与「复化耦合的指数化」无关——后者也常记作 $q$，等于 $\ee^{2\pi\ii\tau}$。事实上，矢量多重态的拉格朗日量具有 $\sim\tau\delta_\epsilon O$ 的形式，因此该背景上的配分函数与复化规范耦合 $\tau$ 无关。

% [Tutor Note: p.23, 这一段里同一个字母 $q$ 出现了三种含义，是全文最易读错的地方之一：
%  (i) 三参数 $(p,q,t)$ 中的 $q$；(ii) 超共形指标变量 $q=\ee^{-\beta}$；(iii) 复化耦合的指数 $q=\ee^{2\pi\ii\tau}$。
%  「$p=0,\ q=t$」说的是把 (i) 中的 $q$ 取成 $t$，紧接着的「$q=\ee^{-\beta}$」却是指 (ii)。
%  原文自己在下一段与 p.24 反复提醒 (ii)(iii) 不同，读者须按上下文区分同一字母的不同含义。]

为行文简洁，我们只使用 $p=0,\ q=t$ 的一维切片。此时 $q=\ee^{-\beta}$。当所讨论的 $\cN=2$ 理论具有味对称性 $G$ 时，我们可以选取元素 $g\in G$，使得绕 $S^1$ 存在由 $g$ 给出的背景 $G$ 规范场。于是该背景上的配分函数就只是 $q=\ee^{-\beta}$ 与 $g\in G$ 的函数。

把 $S^1$ 方向当作时间方向，配分函数可以写成
\begin{equation}
Z(q,g)=\tr_{\cH}(-1)^F q^{\Delta-R}g.
\tag{4.7}
\end{equation}

% [Tutor Note: p.23, (4.7) 中的 $R$ 是 $\SU(2)_R$ 荷的第三个分量（原文 p.23 末段明确说了
%  "the third component of the SU(2)_R charge"），不是 $\U(1)_R$ 荷，也不是别的 $\SU(2)$ 指标；
%  $\Delta$ 是算符的标度维数，$(-1)^F$ 中的 $F$ 是费米子数。态–算符对应把 $S^3$ 上的
%  Hilbert 空间 $\cH$ 与平坦空间中的点状算符一一对应，所以 (4.7) 既可以按「在平坦空间里数算符」
%  来算，也可以按 $S^1\times S^3$ 上的路径积分来算（p.24 开头正是这么说的）。
%  易错点：$q$ 的指数是 $\Delta-R$，不要把 $-R$ 漏掉写成 $\Delta$。]

这里 $\cH$ 是 $S^3$ 上超共形理论的 Hilbert 空间，通过态–算符对应，它等价于平坦空间中共形理论的点状算符空间。$\Delta$ 是算符的标度维数，$R$ 是 $\SU(2)_R$ 荷的第三个分量。由于这种结构，该配分函数常被称为超共形指标。对四维量子场论，这一概念是由 [7,8] 引入的。
